\begin{afterword}
\summaryhead

The post offers the reduction of heat to motion as a simpler example than anger. Heat was once explained by a caloric fluid, and the kinetic theory was controversial. Someone who, like Sadi Carnot in 1824, knows a great deal about heat and about mechanics as separate subjects will feel that heat and motion are independent properties of matter.

Told by a physicist from the future that ``Where there is heat, there is motion, and vice versa,'' such a person might read it as a bridging law between two things, or as saying that they are one thing. Thinker 1 objects that the two words have different meanings; the author answers that different words, like ``2 + 2'' and ``4'', can have one referent, which is found by looking at the world, as with the morning and evening star. Thinker 2 objects that a world where heat and motion come apart is imaginable; the author answers that the physicist stated a rule of inference, not a law, and that what is conceivable is not what is logically possible. Putnam is quoted making that distinction for water and H2O, though the author thinks Putnam uses ``water'' in two senses. The post concludes that reductionism is easy and reduction is hard.

\responsehead

The post's account of why a reduction can seem impossible before it is understood is sound, and its history of heat is right in outline. Kinetic theories were rejected as late as the 1840s, and the facts given about Carnot are correct. Carnot's own later notes support the post's thesis: by 1832 Carnot had written that heat ``is a movement among the particles of bodies.''

Much of the argument had been made before, and the post says as much (``I believe the above has already been said as well?''). It credits Putnam. It does not mention Kripke, who in 1972 used this very example, heat and molecular motion, to explain why an identity discovered by science can seem as if it might have been otherwise. The answer to Thinker 1 is Frege's distinction, with Frege's morning and evening star. Francis Bacon, in 1620, separated the claim that heat generates motion from the claim that heat is motion, the post's two readings. None of this weakens the argument; it shows that the argument is the standard one.

The post's one correction of a philosopher is not supported by the source. It says, with ``seems'', that Putnam takes ``water'' to be ``automatically redefined'' once water is found to be H2O, and says one could hold instead that ``water'' means ``whatever is really in that bottle next to me.'' Putnam's paper rejects the first point: what ``water'' applied to was H2O in 1750, before anyone knew it, and the meaning did not change with the discovery. The alternative is close to Putnam's own account, in which ``this liquid is called water'' makes water whatever is the same liquid as the stuff in the glass. The two readings the post does find in Putnam's two paragraphs correspond to a distinction that a later theory of meaning, two-dimensional semantics, made precise.

A smaller point: heat is not kinetic energy in general but the disordered motion of molecules, and not only their kinetic energy. The argument still holds after the correction.

In short: the standard argument about conceivability, on Kripke's example, with an accurate history of heat, and a reading of Putnam that Putnam's paper does not support; the alternative the post mentions is close to Putnam's own account.
\end{afterword}
